\section{Background}\label{sec:background}

This section fixes notation and collects the results from the literature that are used later.
\S\S\ref{ss:pillowcase}--\ref{ss:decomp} recall the Lagrangian Floer theory of Hedden, Herald and Kirk in
the pillowcase~\cite{HHK2}, which is the setting of Sections~\ref{sec:sheared} and~\ref{sec:twisted}; results
of~\cite{HHK2} are cited by their numbering in arXiv:1501.00028v1. \S\ref{ss:knots} fixes the conventions for
the knots. We use $\SU$ as the group of unit quaternions, so that a
traceless element is a purely imaginary unit quaternion. Coefficients are in $\F_2$ unless stated otherwise.

\subsection{The pillowcase}\label{ss:pillowcase}
Let $G\cong\Z^2\rtimes\Z/2$ be the group of isometries of $\R^2$ generated by
$(\gamma,\theta)\mapsto(\gamma+2\pi,\theta)$, $(\gamma,\theta)\mapsto(\gamma,\theta+2\pi)$ and
$(\gamma,\theta)\mapsto(-\gamma,-\theta)$. The \emph{pillowcase} is the quotient $P=\R^2/G$, a topological
$2$-sphere. The quotient map $\R^2\to P$ is a branched covering, branched over the images of the lattice points
$(\pi\Z)^2$, which are the four \emph{corners} of $P$~\cite[\S3.1]{HHK2}. The rectangle $[0,\pi]\times[0,2\pi]$ is
a fundamental domain; we call $[0,\pi]\times[0,\pi]$ the front face, $[0,\pi]\times[\pi,2\pi]$ the back face, and
the images of the lines $\gamma\in\pi\Z$ the left and right edges. Put $\Pst=(\R^2\setminus(\pi\Z)^2)/G$, oriented
by $d\gamma\wedge d\theta$. The restriction $\R^2\setminus(\pi\Z)^2\to\Pst$ is a regular covering with deck group
$G$, so paths, loops and immersed discs in $\Pst$ lift to $\R^2\setminus(\pi\Z)^2$; we work with such lifts
throughout. The double cover $\R^2/(2\pi\Z)^2$ of $P$ is a torus, and $P$ is its quotient by the elliptic
involution $(\gamma,\theta)\mapsto(-\gamma,-\theta)$.

The pillowcase is the traceless character variety of the four-punctured sphere. With the presentation
$\pi_1(S^2\setminus\{a,b,c,d\})=\langle a,b,c,d\mid ba=cd\rangle$ of~\cite{HHK1,HHK2}, the assignment
\begin{equation}\label{eq:hhk-coords}
a\mapsto\mathbf i,\qquad b\mapsto e^{\gamma\mathbf k}\mathbf i,\qquad c\mapsto e^{\theta\mathbf k}\mathbf i,\qquad
d\mapsto e^{(\theta-\gamma)\mathbf k}\mathbf i
\end{equation}
induces a homeomorphism of $P$ with the space $R(S^2,\{a,b,c,d\})$ of conjugacy classes of $\SU$ representations
sending each of $a,b,c,d$ to a traceless element; the corners correspond to the abelian
representations~\cite[Prop.~3.1]{HHK1},~\cite[Prop.~6.1]{HHK2}. We use these coordinates throughout.

On $\R^2$ we use the function
\[
\varphi=\theta-\gamma .
\]
Translations in $G$ change $\varphi$ by elements of $2\pi\Z$, and $(\gamma,\theta)\mapsto(-\gamma,-\theta)$ changes
its sign. The \emph{diagonal arc} $\Delta=\{(\gamma,\gamma)\mid\gamma\in[0,\pi]\}\subset P$~\cite[(12)]{HHK2} joins
the corners $(0,0)$ and $(\pi,\pi)$, and its preimage in $\R^2$ is the union of the lines
\[
\mathcal D_k=\{\varphi=2\pi k\},\qquad k\in\Z .
\]
The lines $\{\varphi=\pi+2\pi k\}$ are lifts of the arc of slope one on the back face, which is not part of
$\Delta$. The lattice points on $\mathcal D_k$ are the points $(j\pi,j\pi+2\pi k)$, $j\in\Z$.

\subsection{The earring curve}\label{ss:earring}
For $\epsilon>0$ let $L_0=L_0^{\epsilon,0}\colon S^1\to\Pst$ be the immersion defined by
\begin{equation}\label{eq:L0}
\tilde L_0(t)=\bigl(t+\epsilon\sin t+\tfrac\pi2,\ t-\epsilon\sin t+\tfrac\pi2\bigr),\qquad t\in[0,2\pi],
\end{equation}
followed by the branched covering~\cite[Def.~3.4, (13)]{HHK2}; as in~\cite[\S11]{HHK2} we take the perturbation
function $g$ of~\cite[(13)]{HHK2} to be zero. It is the restriction to the boundary of the traceless character
variety of the trivial tangle with an earring, after a small holonomy perturbation~\cite[Thm.~6.2]{HHK2}. For
small $\epsilon$ it is self-transverse with one double point, at $(\pi/2,\pi/2)$, it lies in an
$O(\epsilon)$-neighborhood of $\Delta$, and it passes around the corners $(0,0)$ and $(\pi,\pi)$ through the back
face.

Since $\tilde L_0(t+2\pi)=\tilde L_0(t)+(2\pi,2\pi)$, the preimage of $L_0$ in $\R^2$ is the union of the curves
\[
S^k_+=\tilde L_0+(0,2\pi k),\qquad S^k_-=-\tilde L_0+(0,2\pi k),\qquad k\in\Z,
\]
which we call the \emph{lifts} of $L_0$. Each is a properly embedded line within $\sqrt2\,\epsilon$ of $\mathcal D_k$,
and a translation by $(2\pi a,2\pi b)$ carries $S^k_\pm$ to $S^{k+b-a}_\pm$. Reparametrizing,
$S^0_-=\{(t+\frac\pi2-\epsilon\sin t,\ t+\frac\pi2+\epsilon\sin t)\}$. The slope of $S^k_+$ at the point with parameter $t$ is $(1-\epsilon\cos t)/(1+\epsilon\cos t)$, and that of $S^k_-$ is its reciprocal. Hence on the front
strip $0<\gamma<\pi$ the lift $S^k_+$ has slope less than $1$, and $S^k_-$ has slope greater than $1$ for
$\epsilon<\gamma<\pi-\epsilon$ (it has slope $1$ at $\gamma=\epsilon$ and $\gamma=\pi-\epsilon$); in $P$ these
are the two branches of $L_0$ along the interior of $\Delta$~\cite[\S5.1]{HHK2}. The two lifts near $\mathcal D_k$
cross exactly at the points $(\frac\pi2+j\pi,\frac\pi2+j\pi+2\pi k)$, and they pass each lattice point of $\mathcal D_k$
on opposite sides at distance $\sqrt2\,\epsilon$: for instance $S^0_+$ passes $(0,0)$ through $(-\epsilon,\epsilon)$
and $S^0_-$ through $(\epsilon,-\epsilon)$, and the sides alternate from one lattice point of $\mathcal D_k$ to the next.

\subsection{Restricted curves and the Floer complex}\label{ss:restricted}
Let $\ell_1$ be the line field of slope one on $\R^2$; it descends to $\Pst$. For a path $\alpha$ immersed in $\Pst$,
the Maslov index $\mu(\alpha,\ell_1)$ is the signed number of times the tangent line of $\alpha$ passes $\ell_1$,
counterclockwise passages counting $+1$, with the half-open endpoint convention of~\cite[\S2.3]{HHK2}. For lines
$\ell_0,\ell_1',\ell$, pairwise transverse, the \emph{triple index} $\tau(\ell_0,\ell_1',\ell)\in\{0,1\}$ is $1$
exactly when the clockwise rotation of $\ell_0$ to $\ell_1'$ passes through $\ell$~\cite[Def.~2.5]{HHK2}. Let $z\in
H^1(\Pst;\Z/4)$ be the class taking the value $1$ on a small counterclockwise loop around a
corner~\cite[Def.~3.1]{HHK2}.

An immersed circle in $\Pst$ is \emph{unobstructed} if every lift of it to the universal cover of $\Pst$ is a
properly embedded line; by a lemma of Abouzaid~\cite[Lemma~2.2]{Abouzaid} this holds if and only if it is
homotopically essential and contains no \emph{fishtail}, that is, no subinterval whose endpoints map to the same
point and which maps to a nullhomotopic loop~\cite[Def.~2.1]{HHK2}. A \emph{proper immersion} of an interval is
the image of a smooth immersion $I\to\R^2$ sending the endpoints to lattice points and the interior into
$\R^2\setminus(\pi\Z)^2$; its slopes at the endpoints are its \emph{limiting slopes}~\cite[Def.~3.3]{HHK2}.

\begin{definition}[{\cite[Def.~3.6 and p.~33]{HHK2}}]\label{def:restricted}
A \emph{restricted immersed circle} is an unobstructed immersed circle $L\colon S^1\to\Pst$ with
$\mu(L,\ell_1)+z(L)\equiv0\pmod4$. A \emph{restricted immersed arc} is a proper immersion of an interval whose
restriction to the complement of small discs around the corners contains no fishtail. A \emph{restricted immersed
$1$-manifold} is a disjoint union of one restricted immersed arc and finitely many restricted immersed circles.
\end{definition}

The earring curve $L_0$ is unobstructed with $\mu(L_0,\ell_1)=0=z(L_0)$, so it is a restricted immersed
circle~\cite[pp.~16,~18]{HHK2}. The image of a straight segment joining two lattice points whose interior misses the
lattice is a restricted immersed arc~\cite[p.~18]{HHK2}.

\begin{definition}[{\cite[Def.~3.11]{HHK2}}]\label{def:admissible}
A pair $(L_0,L_1)$ of restricted immersed curves is \emph{admissible} if at least one of them is a circle, if $L_0$
and $L_1$ intersect transversely, and if, whenever loops $\alpha_0$ in the domain of $L_0$ and $\alpha_1$ in the
domain of $L_1$ have $L_0\circ\alpha_0$ freely homotopic to $L_1\circ\alpha_1$ in $\Pst$, both $\alpha_0$ and
$\alpha_1$ are nullhomotopic.
\end{definition}

For an admissible pair, $C(L_0,L_1)$ is the $\F_2$-vector space on $L_0\cap L_1$. A \emph{bigon} from $p$ to $p'$ is
an orientation-preserving immersion of a disc into $\Pst$ with two convex corners, at $p$ and $p'$, whose
boundary, traversed counterclockwise, runs along $L_0$ from $p$ to $p'$ and then along $L_1$ from $p'$ back to $p$;
equivalent bigons are identified~\cite[Def.~2.10, (5)]{HHK2}. The differential $\partial p=\sum_{p'}\#\mathcal
M(p,p')\,p'$ counts bigons from $p$ to $p'$ modulo $2$, and by~\cite[Cor.~3.13]{HHK2} each $\mathcal M(p,p')$ contains
at most one bigon. A bigon lifts to an immersed disc in $\R^2\setminus(\pi\Z)^2$, so its two corners lie on the
same component of $L_1$; there is no differential between generators on different
components~\cite[p.~37]{HHK2}. We write $\HF(L_0,L_1)$ for the homology.

The complex carries a relative $\Z/4$-grading $\gr(p,p')=\gr(p)-\gr(p')$, defined for generators on the same
component of $L_1$~\cite[Def.~3.7]{HHK2}. The differential lowers it by one, and a bigon from $p$ to $p'$ forces
$\gr(p,p')=1$~\cite[Prop.~3.8]{HHK2}. If $\alpha_0$ is a path in $L_0$ from $p$ to $p'$ and $\alpha_1$ a path in $L_1$
from $p'$ to $p$, then~\cite[Lemma~3.10]{HHK2}
\begin{equation}\label{eq:grading15}
\gr(p,p')=\mu(\alpha_0,\ell_1)+\mu(\alpha_1,\ell_1)+\tau(T_pL_0,T_pL_1,\ell_1)-\tau(T_{p'}L_0,T_{p'}L_1,\ell_1)
+z(\alpha_0*\alpha_1)\pmod4 .
\end{equation}
Two consequences are recorded in~\cite[\S\S5.1--5.2]{HHK2}.
\begin{enumerate}
\item If $L_1$ meets the interior of $\Delta$ transversely at $x$, then for small $\epsilon$ there are exactly two
generators near $x$: $x^+$ on the branch of $L_0$ of slope less than $1$ and $x^-$ on the branch of slope greater
than $1$. They satisfy $\gr(x^+,x^-)=1$.
\item If $p_+,p'_+$ lie on the branch of slope less than $1$ on the front face, $\alpha_0$ runs along that branch, and
$L_1$ is transverse to $\Delta$ at both points, then
\begin{equation}\label{eq:grading16}
\gr(p_+,p'_+)=\mu(\alpha_1,\ell_1)+z(\alpha_0*\alpha_1)\pmod4 ,
\end{equation}
and the same formula holds for two points on the branch of slope greater than $1$.
\end{enumerate}
In terms of lifts, $x^\pm$ are the intersections of a lift of $L_1$ with the lifts $S^k_\pm$ of $L_0$ near a point of
$\mathcal D_k$ on the front strip.

Suppose $L_1$ is a restricted immersed arc with an endpoint at the corner $(0,0)$ whose limiting slope is bounded
away from $1$. Then for small $\epsilon$ there is a unique intersection point $r_+$ of $L_0$ and $L_1$ converging
to that corner as $\epsilon\to0$~\cite[Lemma~5.1]{HHK2}. When $L_1$ comes from a knot $K$ as in \S\ref{ss:decomp},
the grading of the component containing the arc is made absolute by
\begin{equation}\label{eq:grading19}
\gr(r_+)=\sigma(K)\pmod4
\end{equation}
\cite[(19)]{HHK2}. Since $\sigma(K)$ is even, \eqref{eq:grading19} does not depend on the sign convention for the
signature. On circle components Hedden, Herald and Kirk fix the absolute grading by matching one generator with
the corresponding generator of Kronheimer and Mrowka's complex~\cite[\S6.1]{HHK2},~\cite[Prop.~4.4]{KM2}.

\begin{theorem}[{\cite[Thm.~5.4]{HHK2}}]\label{thm:hhk54}
Let $L_1\colon S^1\to\Pst$ be a restricted immersed circle that misses the left and right edges of the
pillowcase, and let $\tilde L_1\colon[0,2\pi]\to\R^2$ be a lift of $t\mapsto L_1(e^{it})$. Then its vertical degree
$d=\frac1{2\pi}|\theta(\tilde L_1(2\pi))-\theta(\tilde L_1(0))|$ is even, and $\HF(L_0,L_1)$ has rank $d/2$ in each of
the four gradings.
\end{theorem}

\begin{remark}\label{rem:invariance}
By~\cite[Thm.~4.1, Cor.~4.6]{HHK2}, $\HF(L_0,L_1)$ is independent of $\epsilon$ and, as a relatively graded group,
depends only on the free homotopy classes of the circle components of $L_1$ and the homotopy class rel endpoints
of its arc. The computations in Section~\ref{sec:sheared} are direct and do not use this invariance.
Section~\ref{sec:twisted} uses it only through Theorem~\ref{thm:hhk54}, for the circle components: in~\cite{HHK2}
that theorem is deduced from~\cite[Cor.~4.6]{HHK2}, which follows from~\cite[Thm.~4.1]{HHK2}.
\end{remark}

\subsection{Tangle decompositions and the conjecture of Hedden, Herald and Kirk}\label{ss:decomp}
Let $K\subset S^3$ be a knot, decomposed along a Conway sphere as $(S^3,K)=(Y,T)\cup(D,U)$, where $(D,U)$ is the
trivial $2$-tangle in a ball and the common boundary is identified with
$(S^2,\{a,b,c,d\})$~\cite[(17)]{HHK2}. Let $R(Y,T)$ be the space of conjugacy classes of traceless $\SU$
representations of $\pi_1(Y\setminus T)$, and $R_\pi(Y,T)$ its version perturbed by a holonomy perturbation $\pi$
supported in the interior of $Y$. Restriction to the boundary gives $L_1\colon R_\pi(Y,T)\to P$, and the earring
curve $L_0$ of \S\ref{ss:earring} comes from the trivial side. When $L_1$ is a restricted immersed $1$-manifold and
$(L_0,L_1)$ is admissible, Hedden, Herald and Kirk define
\[
\Hnat(Y,T,\pi)=\HF(L_0,L_1),
\]
which is the direct sum of the contributions of the components of
$R_\pi(Y,T)$~\cite[p.~34, \S6.1]{HHK2}. The two abelian traceless representations of $\pi_1(Y\setminus T)$ are the
endpoints of the arc component; the one that extends over the knot complement maps to the corner $(0,0)$, and its
perturbation gives the generator $r_+$~\cite[\S6.1, \S8]{HHK2}. The remaining generators come in pairs $x^\pm$, one
pair for each transverse intersection point of $L_1$ with the interior of $\Delta$. For a torus knot, with the
tangle side unperturbed and a small perturbation on the side of the earring, these intersection points correspond
to the irreducible traceless representations of the knot group, and there are $1+|\sigma(K)|$
generators~\cite[Thm.~11.1]{HHK1},~\cite[\S11.10]{HHK2}.

\begin{conjecture}[{\cite[Conj.~6.5]{HHK2}}]\label{conj:hhk65}
``Given a knot $(X,K)$ in a homology 3-sphere, there exists a 2-tangle decomposition as in Equation (17), such that
for suitably small generic perturbations $\pi$, $L_1:R_\pi(Y,T)\to P$ is a restricted immersed 1-manifold and
$H^\natural(Y,T,\pi)$ is isomorphic to the reduced instanton homology $I^\natural(X,K)$.''
\end{conjecture}

Their Equation (17) is the decomposition $(X,K)=(Y,T)\cup(D,U)$ along a Conway sphere, which for $X=S^3$ is the
decomposition $(S^3,K)=(Y,T)\cup(D,U)$ of this section.

The conjecture asserts the existence of one good decomposition, and the answer does depend on the decomposition:
different choices of the gluing data give different curves $L_1$~\cite[p.~55]{HHK2}, and with the earring on the
other tangle the pillowcase homology of $T(4,5)$ has the wrong rank~\cite[\S10.3]{CHKK}. Hedden, Herald and Kirk
verify the conjecture for all two-bridge knots~\cite[\S7]{HHK2}. For the torus knots $T(3,n)$ with $n\ge4$ prime
to $3$, suitable decompositions from~\cite[\S\S10--11]{HHK2} satisfy it, with the $\Z/4$-gradings~\cite{WuebbenV}. Section~\ref{sec:sheared}
obtains the decompositions used for $P(-2,3,q)$ from these by regluing.

\subsection{The knots}\label{ss:knots}
We use the convention that positive knots have negative signature, as do~\cite{DS19,DS24,WuebbenV} and the tables
of~\cite{HHK2}; for instance $\sigma(T(3,5))=-8$~\cite[p.~60]{HHK2}. Braid generators are written $\sigma_1,\sigma_2$
and are positive crossings, so that the closure of $(\sigma_1\sigma_2)^n$ is the positive torus knot $T(3,n)$.

For odd $q\ge1$, $K_q=P(-2,3,q)$ is the pretzel knot with three twist regions of $-2$, $3$ and $q$ half-twists, with
the chirality for which its first three members are the positive torus knots $P(-2,3,1)=T(2,5)$,
$P(-2,3,3)=T(3,4)$ and $P(-2,3,5)=T(3,5)$; for $q\ge7$ it is hyperbolic (Proposition~\ref{prop:hyperbolic}). Section~\ref{sec:sheared}
decomposes it by regluing Hedden, Herald and Kirk's decomposition of $T(3,5)$ (Theorem~\ref{thm:topology}).

For $n\ge7$ with $n\equiv1,2\pmod 6$ and $m\ge1$, the twisted torus knot $T(3,n;2,m)$ is the closure of the
three-strand braid
\[
(\sigma_1\sigma_2)^n\,\sigma_1^{2m},
\]
that is, $T(3,n)$ with $m$ positive full twists added on two adjacent strands. It is a positive braid closure.
Section~\ref{sec:twisted} treats these knots with decompositions obtained from those of $T(3,n)$ in the same way.
